\documentclass[11pt]{ctexart}
\usepackage[a4paper,margin=1in]{geometry}
\usepackage{longtable,booktabs}
\title{Geant4-Agent 回归测试报告（中文）}
\author{自动化评测流程}
\date{2026-03-06\_124519}
\begin{document}
\maketitle
\section{统一 Pass/Fail 总表}
\begin{longtable}{p{0.06\linewidth}p{0.08\linewidth}p{0.07\linewidth}p{0.08\linewidth}p{0.07\linewidth}p{0.07\linewidth}p{0.07\linewidth}p{0.07\linewidth}p{0.09\linewidth}p{0.06\linewidth}p{0.21\linewidth}}
\toprule
ID & 领域 & 风格 & 参数完整性 & 期望初始 & 实际初始 & 期望最终 & 实际最终 & 缺参(初->终) & 结果 & 失败原因 \\
\midrule
MTM01 & medical & fuzzy & missing & 失败 & 失败 & 通过 & 通过 & 13->0 & 通过 & - \\
MTM02 & aerospace & precise & missing & 失败 & 失败 & 通过 & 通过 & 2->0 & 通过 & - \\
MTM03 & industrial\_ndt & precise & missing & 失败 & 失败 & 通过 & 通过 & 8->0 & 通过 & - \\
MTM04 & physics & precise & missing & 失败 & 失败 & 通过 & 通过 & 8->0 & 通过 & - \\
MTM05 & nuclear\_engineering & fuzzy & missing & 失败 & 失败 & 通过 & 通过 & 13->0 & 通过 & - \\
MTM06 & security\_screening & fuzzy & missing & 失败 & 失败 & 通过 & 通过 & 4->0 & 通过 & - \\
MTM07 & semiconductor & precise & missing & 失败 & 失败 & 通过 & 通过 & 13->0 & 通过 & - \\
MTM08 & education & fuzzy & missing & 失败 & 失败 & 通过 & 通过 & 11->0 & 通过 & - \\
MTM09 & environmental\_radiation & fuzzy & missing & 失败 & 失败 & 通过 & 通过 & 13->0 & 通过 & - \\
MTM10 & high\_energy\_physics & precise & missing & 失败 & 失败 & 通过 & 通过 & 10->0 & 通过 & - \\
\bottomrule
\end{longtable}
\section{基准对比}
\begin{itemize}
\item pass\_rate=1.000
\item initial\_alignment\_rate=1.000
\item final\_alignment\_rate=1.000
\item missing\_reduction\_rate=1.000
\item confirm\_apply\_success\_rate=1.000
\item rollback\_after\_confirm\_rate=0.000
\item delta\_pass\_rate=+0.050
\item regressions=无
\item improvements=无
\end{itemize}
\appendix
\section{附录：完整对话与内部处理链（自然语言）}
\subsection{案例 MTM01（medical）}
\textbf{对话：}
\begin{itemize}
\item 用户：???????????????????
\item 系统：你希望构建哪种几何结构呢？比如一个简单的长方体，还是一个圆柱体？另外，长方体的X方向尺寸大概是多少？
\item 用户：?????? 300 mm x 300 mm x 300 mm??? G4\_WATER?source.type=beam?particle=gamma?energy=6 MeV?position=(0,0,-400) mm?direction=+z?physics=QBBC?output=root?
\item 系统：好的，配置已完成！根据你提供的信息，我已经为你设置好了几何结构（300mm x 300mm x 300mm的长方体）、材料（G4\_WATER）、粒子源（6 MeV伽马射线束）、物理列表（QBBC）和输出格式（ROOT文件）。
\end{itemize}
\textbf{内部链路（自然语言）：}
\begin{itemize}
\item 第 1 轮：动作=ask\_clarification，阶段=geometry，闭环状态=False。
\item 缺失字段：geometry.structure、geometry.params.module\_x、geometry.params.module\_y、geometry.params.module\_z、materials.volume\_material\_map、source.type、source.particle、source.energy、source.position、source.direction、physics.physics\_list、output.format、output.path。
\item LLM 使用=True；回退=无；后端=llm\_semantic\_frame+runtime\_semantic。
\item 第 2 轮：动作=finalize，阶段=materials，闭环状态=True。
\item 缺失字段：无。
\item LLM 使用=True；回退=无；后端=llm\_slot\_frame+runtime\_semantic。
\end{itemize}
\subsection{案例 MTM02（aerospace）}
\textbf{对话：}
\begin{itemize}
\item 用户：I need a first-pass proton irradiation setup for a silicon device; keep it simple for now.
\item 系统：Here's a summary of your setup so far:

**Geometry:** Updated with a simple structure and module dimensions.
**Materials:** Updated. Silicon is selected for your device.
**Source:** Updated for proton irradiation.
**Physics:** Physics list has been selected.

**Next, to run the simulation, please specify:**
*   The output format (e.g., ROOT, CSV).
*   The output file path.

Once these are provided, your configuration will be ready.
\item 用户：Use a silicon box 50 mm x 50 mm x 5 mm, beam source proton 100 MeV from (0,0,-100) mm along +z, physics FTFP\_BERT, output json.
\item 系统：Your new geometry dimensions would overwrite previously confirmed values. Please confirm this change:

*   `module\_x`: 1000.0 mm → 50.0 mm
*   `module\_y`: 1000.0 mm → 50.0 mm
\item 用户：confirm
\item 系统：You are about to overwrite previously confirmed geometry parameters. Please confirm this change:

*   `module\_x`: 1000.0 mm → 50.0 mm
*   `module\_y`: 1000.0 mm → 50.0 mm
\end{itemize}
\textbf{内部链路（自然语言）：}
\begin{itemize}
\item 第 1 轮：动作=summarize\_progress，阶段=materials，闭环状态=False。
\item 缺失字段：output.format、output.path。
\item LLM 使用=True；回退=无；后端=llm\_slot\_frame+runtime\_semantic。
\item 第 2 轮：动作=confirm\_overwrite，阶段=source，闭环状态=False。
\item 缺失字段：无。
\item LLM 使用=True；回退=无；后端=llm\_slot\_frame+runtime\_semantic。
\item 第 3 轮：动作=confirm\_overwrite，阶段=physics，闭环状态=True。
\item 缺失字段：无。
\item LLM 使用=True；回退=无；后端=llm\_semantic\_frame+runtime\_semantic+confirmed\_pending\_overwrite。
\end{itemize}
\subsection{案例 MTM03（industrial\_ndt）}
\textbf{对话：}
\begin{itemize}
\item 用户：???????????????? gamma ???
\item 系统：当前配置进度如下：

**几何**：已更新结构类型和尺寸参数（module\_x、module\_y、module\_z）。  
**材料**：还需设置体积与材料的对应关系。  
**源**：已选定粒子类型（gamma），并记录了选择来源和依据；还需补充源的具体类型、能量、位置和方向。  
**物理**：待选择物理列表。  
**输出**：待指定输出格式和保存路径。

请继续补充剩余参数，完成后我会为你生成完整的配置文件。
\item 用户：??????????? 40 mm??? 80 mm??? G4\_Fe?source.type=point?particle=gamma?energy=2 MeV?position=(0,0,-200) mm?direction=+z?physics=Shielding?output=csv?
\item 系统：检测到需要覆盖已确认的配置项，请确认是否应用以下修改：

- 几何结构类型：从 single\_box 改为 single\_tubs
- 几何参数 child\_rmax：从空值设为 40.0
\item 用户：确认
\item 系统：好的，检测到您需要修改已确认的几何配置。请确认是否应用以下更改：

*   几何结构类型：从 **single\_box** 改为 **single\_tubs**
*   几何参数 child\_rmax：从 **空值** 设为 **40.0**
\end{itemize}
\textbf{内部链路（自然语言）：}
\begin{itemize}
\item 第 1 轮：动作=summarize\_progress，阶段=materials，闭环状态=False。
\item 缺失字段：materials.volume\_material\_map、source.type、source.energy、source.position、source.direction、physics.physics\_list、output.format、output.path。
\item LLM 使用=True；回退=无；后端=llm\_slot\_frame+runtime\_semantic。
\item 第 2 轮：动作=confirm\_overwrite，阶段=source，闭环状态=False。
\item 缺失字段：无。
\item LLM 使用=True；回退=无；后端=llm\_slot\_frame+runtime\_semantic。
\item 第 3 轮：动作=confirm\_overwrite，阶段=physics，闭环状态=True。
\item 缺失字段：无。
\item LLM 使用=True；回退=无；后端=llm\_semantic\_frame+runtime\_semantic+confirmed\_pending\_overwrite。
\end{itemize}
\subsection{案例 MTM04（physics）}
\textbf{对话：}
\begin{itemize}
\item 用户：Gamma point source 1 MeV at (0,0,-100) toward +z; please build the rest.
\item 系统：Here's your current setup:

✅ **Source** is configured: gamma point source, 1 MeV, positioned at (0,0,-100) and directed along +z.

❌ **Geometry** is incomplete. You still need to define the structure and its parameters (module\_x, module\_y, module\_z).
❌ **Materials** are incomplete. You still need to assign materials to your geometry volumes.
❌ **Physics** list is not yet selected.
❌ **Output** format and file path are not yet specified.

Please provide the missing geometry details to continue.
\item 用户：Set copper box 1 m x 1 m x 1 m, physics FTFP\_BERT, output hdf5.
\item 系统：Your configuration is complete.

✅ **Source**: Gamma point source, 1 MeV, at (0,0,-100) directed along +z.
✅ **Geometry**: Copper box, 1 m x 1 m x 1 m.
✅ **Physics**: FTFP\_BERT physics list.
✅ **Output**: HDF5 format.
✅ **Materials**: Copper assigned to the geometry volume.
\end{itemize}
\textbf{内部链路（自然语言）：}
\begin{itemize}
\item 第 1 轮：动作=summarize\_progress，阶段=geometry，闭环状态=False。
\item 缺失字段：geometry.structure、geometry.params.module\_x、geometry.params.module\_y、geometry.params.module\_z、materials.volume\_material\_map、physics.physics\_list、output.format、output.path。
\item LLM 使用=True；回退=无；后端=llm\_slot\_frame+runtime\_semantic。
\item 第 2 轮：动作=finalize，阶段=materials，闭环状态=True。
\item 缺失字段：无。
\item LLM 使用=True；回退=无；后端=llm\_slot\_frame+runtime\_semantic。
\end{itemize}
\subsection{案例 MTM05（nuclear\_engineering）}
\textbf{对话：}
\begin{itemize}
\item 用户：?????????????????????
\item 系统：你打算构建哪种几何结构呢？比如一个简单的盒子，还是圆柱体？另外，盒子的X方向尺寸（module\_x）需要设定为多少？
\item 用户：????? 200 mm x 200 mm x 200 mm??? G4\_Pb?source.type=point?particle=neutron?energy=14 MeV?position=(0,0,-50) mm?direction=+z?physics=QGSP\_BERT\_HP?output=xml?
\item 系统：配置已完成。
\end{itemize}
\textbf{内部链路（自然语言）：}
\begin{itemize}
\item 第 1 轮：动作=ask\_clarification，阶段=geometry，闭环状态=False。
\item 缺失字段：geometry.structure、geometry.params.module\_x、geometry.params.module\_y、geometry.params.module\_z、materials.volume\_material\_map、source.type、source.particle、source.energy、source.position、source.direction、physics.physics\_list、output.format、output.path。
\item LLM 使用=True；回退=无；后端=llm\_semantic\_frame+runtime\_semantic。
\item 第 2 轮：动作=finalize，阶段=materials，闭环状态=True。
\item 缺失字段：无。
\item LLM 使用=True；回退=无；后端=llm\_slot\_frame+runtime\_semantic。
\end{itemize}
\subsection{案例 MTM06（security\_screening）}
\textbf{对话：}
\begin{itemize}
\item 用户：Need a quick cargo-screening toy setup with forward photons; details can follow.
\item 系统：Here's your current setup status:

**Updated:**
- **Geometry:** Structure and module dimensions (X, Y) are set.
- **Source:** Type, particle (photon), and energy are configured.

**Still needed:**
- **Materials:** Assign materials to your volumes.
- **Physics:** Select a physics list.
- **Output:** Choose a format and specify a save location.

Please provide the missing information to complete the configuration.
\item 用户：Use steel cylinder target radius 60 mm half-length 120 mm, beam gamma 3 MeV from (0,0,-500) mm to +z, physics Shielding, output root.
\item 系统：You have existing confirmed settings that would be changed. Please confirm this update:

*   **Geometry structure:** `single\_box` → `single\_tubs`
*   **Source energy:** `1.0` → `3.0`
\item 用户：confirm
\item 系统：You have existing confirmed settings that would be changed. Please confirm this update:

*   **Geometry structure:** `single\_box` → `single\_tubs`
*   **Source energy:** `1.0` → `3.0`
\end{itemize}
\textbf{内部链路（自然语言）：}
\begin{itemize}
\item 第 1 轮：动作=summarize\_progress，阶段=materials，闭环状态=False。
\item 缺失字段：materials.volume\_material\_map、physics.physics\_list、output.format、output.path。
\item LLM 使用=True；回退=无；后端=llm\_slot\_frame+runtime\_semantic。
\item 第 2 轮：动作=confirm\_overwrite，阶段=source，闭环状态=False。
\item 缺失字段：无。
\item LLM 使用=True；回退=无；后端=llm\_slot\_frame+runtime\_semantic。
\item 第 3 轮：动作=confirm\_overwrite，阶段=physics，闭环状态=True。
\item 缺失字段：无。
\item LLM 使用=True；回退=无；后端=llm\_semantic\_frame+runtime\_semantic+confirmed\_pending\_overwrite。
\end{itemize}
\subsection{案例 MTM07（semiconductor）}
\textbf{对话：}
\begin{itemize}
\item 用户：????????????????????????
\item 系统：你希望构建哪种几何结构呢？比如，是一个简单的盒子，还是一个圆柱体？另外，这个结构的X方向尺寸需要设定为多少？
\item 用户：?? 20 mm x 20 mm x 1 mm ? G4\_Si ????source.type=isotropic?particle=neutron?energy=2 MeV?position=(0,0,0) mm?physics=QGSP\_BERT\_HP?output=json?
\item 系统：配置已完成。
\end{itemize}
\textbf{内部链路（自然语言）：}
\begin{itemize}
\item 第 1 轮：动作=ask\_clarification，阶段=geometry，闭环状态=False。
\item 缺失字段：geometry.structure、geometry.params.module\_x、geometry.params.module\_y、geometry.params.module\_z、materials.volume\_material\_map、source.type、source.particle、source.energy、source.position、source.direction、physics.physics\_list、output.format、output.path。
\item LLM 使用=True；回退=无；后端=llm\_semantic\_frame+runtime\_semantic。
\item 第 2 轮：动作=finalize，阶段=materials，闭环状态=True。
\item 缺失字段：无。
\item LLM 使用=True；回退=无；后端=llm\_slot\_frame+runtime\_semantic。
\end{itemize}
\subsection{案例 MTM08（education）}
\textbf{对话：}
\begin{itemize}
\item 用户：For classroom demo, I want gamma attenuation in copper; you can ask only essential missing parts.
\item 系统：Here's your current setup:

**Geometry:** Still needs the structure type and the module X dimension.
**Materials:** Copper is selected. Still needs to assign it to a volume.
**Source:** Gamma particle is selected. Still needs the source type and energy.
**Physics:** Physics list is configured.
**Output:** Still needs the format and file path.

Please provide the missing geometry details first.
\item 用户：Okay, copper box 100 mm x 100 mm x 100 mm, point gamma 1 MeV at (0,0,-200) mm, direction +z, physics FTFP\_BERT, output csv.
\item 系统：You're about to overwrite a confirmed value. Please confirm this change:

**Physics list:** Shielding → FTFP\_BERT
\item 用户：confirm
\item 系统：You've already confirmed the physics list as "Shielding." Changing it to "FTFP\_BERT" will overwrite that. Please confirm this change:

**Physics list:** Shielding → FTFP\_BERT
\end{itemize}
\textbf{内部链路（自然语言）：}
\begin{itemize}
\item 第 1 轮：动作=summarize\_progress，阶段=geometry，闭环状态=False。
\item 缺失字段：geometry.structure、geometry.params.module\_x、geometry.params.module\_y、geometry.params.module\_z、materials.volume\_material\_map、source.type、source.energy、source.position、source.direction、output.format、output.path。
\item LLM 使用=True；回退=无；后端=llm\_slot\_frame+runtime\_semantic。
\item 第 2 轮：动作=confirm\_overwrite，阶段=materials，闭环状态=False。
\item 缺失字段：无。
\item LLM 使用=True；回退=无；后端=llm\_slot\_frame+runtime\_semantic。
\item 第 3 轮：动作=confirm\_overwrite，阶段=source，闭环状态=True。
\item 缺失字段：无。
\item LLM 使用=True；回退=无；后端=llm\_semantic\_frame+runtime\_semantic+confirmed\_pending\_overwrite。
\end{itemize}
\subsection{案例 MTM09（environmental\_radiation）}
\textbf{对话：}
\begin{itemize}
\item 用户：???????????????????????
\item 系统：你希望构建什么形状的几何结构呢？另外，这个结构在X方向的尺寸是多少？
\item 用户：??? 500 mm x 500 mm x 500 mm ?????? G4\_AIR?source.type=point?particle=gamma?energy=0.5 MeV?position=(0,0,-300) mm?direction=+z?physics=QBBC?output=root?
\item 系统：配置已完成。
\end{itemize}
\textbf{内部链路（自然语言）：}
\begin{itemize}
\item 第 1 轮：动作=ask\_clarification，阶段=geometry，闭环状态=False。
\item 缺失字段：geometry.structure、geometry.params.module\_x、geometry.params.module\_y、geometry.params.module\_z、materials.volume\_material\_map、source.type、source.particle、source.energy、source.position、source.direction、physics.physics\_list、output.format、output.path。
\item LLM 使用=True；回退=无；后端=llm\_semantic\_frame+runtime\_semantic。
\item 第 2 轮：动作=finalize，阶段=materials，闭环状态=True。
\item 缺失字段：无。
\item LLM 使用=True；回退=无；后端=llm\_slot\_frame+runtime\_semantic。
\end{itemize}
\subsection{案例 MTM10（high\_energy\_physics）}
\textbf{对话：}
\begin{itemize}
\item 用户：Set geometry as a sphere target; I have not decided radius or source details yet.
\item 系统：Here's your current setup:

**Geometry:** You've set a sphere target and its root volume. Still need the sphere's maximum radius (child\_rmax).

**Materials:** Still need to assign materials to your volumes.

**Source:** Still need the source type, particle type, and particle energy.

**Physics \& Output:** Still need the physics list, output format, and output path.

Let me know what you'd like to set next.
\item 用户：Sphere target radius 80 mm, lead material, electron point source 5 MeV at (0,0,-200) mm toward +z, physics FTFP\_BERT, output hdf5.
\item 系统：Your simulation is ready. You've configured a sphere target with an 80 mm radius, made of lead. The source is a 5 MeV electron point at (0,0,-200) mm, directed along +Z. The physics list is FTFP\_BERT, and the output format is HDF5.
\end{itemize}
\textbf{内部链路（自然语言）：}
\begin{itemize}
\item 第 1 轮：动作=summarize\_progress，阶段=geometry，闭环状态=False。
\item 缺失字段：geometry.params.child\_rmax、materials.volume\_material\_map、source.type、source.particle、source.energy、source.position、source.direction、physics.physics\_list、output.format、output.path。
\item LLM 使用=True；回退=无；后端=llm\_slot\_frame+runtime\_semantic。
\item 第 2 轮：动作=finalize，阶段=materials，闭环状态=True。
\item 缺失字段：无。
\item LLM 使用=True；回退=无；后端=llm\_slot\_frame+runtime\_semantic。
\end{itemize}
\end{document}